\documentclass{lug}

\title{Bottle}
\author{Jack Rosenthal}
\date{2017-03-23}
\institute{Mines Linux Users Group}

\begin{document}

\begin{frame}{HTTP Crash Course}
    \begin{itemize}[<+->]
        \item \texttt{GET} -- Your typical ``downloads from a web page'' HTTP
            request
            \begin{itemize}
                \item Has no support for sending data to the server other than
                    the URL requested
            \end{itemize}
        \item \texttt{POST} -- Your typical ``form submission'' HTTP request
            \begin{itemize}
                \item Sends multipart form data along with the request
                \item Web browsers ask you to ``resubmit'' when refreshed
            \end{itemize}
        \item \texttt{PUT} -- Ask to put a file at a location
            \begin{itemize}
                \item Not needed as the same data can be sent using
                    \texttt{POST}
                \item No support in older browsers
                \item Simple to use when there's no web browser involved
            \end{itemize}
    \end{itemize}
\end{frame}

\begin{frame}{Decorators}
    Decorators are a pretty way to wrap functions using functions that return
    functions.

    \pause

    Both the following are equivalent:
    \inputminted{python3}{examples/decorator_equiv.py}

    Bottle makes heavy use of decorators to bind into routing.
\end{frame}

\begin{frame}{Bottle: Really Simple Web Framework}
    \begin{columns}
        \begin{column}{0.7\linewidth}
            \begin{itemize}[<+->]
                \item Provides routing and convenient access to data
                \item Built in HTTP server, or use any WSGI-compatible web
                    server
                \item Very lightweight, only a couple thousand lines of code
            \end{itemize}
        \end{column}
        \begin{column}{0.3\linewidth}
            \includegraphics[width=\linewidth]{graphics/bottle}
        \end{column}
    \end{columns}

    \pause[\thebeamerpauses]
    \bigskip

    \begin{block}{A \emph{Hello, World!} App}
        \begin{minipage}{\linewidth}
            \tiny
            \inputminted{python3}{examples/bottle_ex.py}
        \end{minipage}
    \end{block}
\end{frame}

\begin{frame}{Where should I use Bottle?}
    \textbf{What Bottle Is}

    \begin{itemize}[<+->]
        \item Bottle is a \emph{micro-framework}
        \item Bottle is only a library
        \item Bottle really small and fast
    \end{itemize}

    \pause[\thebeamerpauses]

    \textbf{What Bottle Is Not}

    \begin{itemize}[<+->]
        \item Bottle is not a MVC framework
        \item Bottle will not generate files for you
        \item \emph{No magic included!}
    \end{itemize}
\end{frame}

\begin{frame}{Routing}
    Routing is how you let Bottle know which URLs map to which functions.
    Bottle uses decorators to specify this.

    \inputminted{python3}{examples/route_simple.py}

    \pause
    \medskip

    You can specify multiple routes per function, even using default arguments.

    \inputminted{python3}{examples/multiroute.py}
\end{frame}

\begin{frame}{More Routing}
    URLs match more specific routes before more generalized.

    \inputminted{python3}{examples/generalized.py}

    \medskip

    So \texttt{/page/WikiPage} will use \texttt{wiki\_page}, but
    \texttt{/jrosenth/home} will use \texttt{user\_page}.
\end{frame}

\begin{frame}{Routing Modifiers}
    You can modify parameters to only match certain types.

    \begin{itemize}[<+->]
        \item \texttt{:int} -- Only match integers
        \item \texttt{:float} -- Only match real numbers
        \item \texttt{:path} -- Match path to end of URL (including slashes)
        \item \texttt{:re} -- Match a regular expression
    \end{itemize}

    \pause[\thebeamerpauses]
    \medskip

    \inputminted{python3}{examples/route_modifiers.py}

\end{frame}

\begin{frame}{Routing Methods}
    The \texttt{route} decorator accepts all HTTP methods by default, you can
    use any of the alternate decorators to accept specific methods.

    \medskip

    \inputminted{python3}{examples/route_methods.py}
\end{frame}

\begin{frame}{Error Routes}
    The \texttt{error} decorator matches certain HTTP errors.

    \medskip

    \inputminted{python3}{examples/route_error.py}
\end{frame}

\begin{frame}{Serving up a static directory}
    \inputminted{python3}{examples/static_files.py}
\end{frame}

\begin{frame}{Accessing \texttt{POST} data}
    The \texttt{request} object gives you access to information about the
    request, including a dictionary contating form data called \texttt{forms}.

    \inputminted{python3}{examples/form_data.py}
\end{frame}

\begin{frame}{Mmm! Cookies!}
    \inputminted{python3}{examples/cookies.py}
\end{frame}

\begin{frame}{So what can you \texttt{return}?}
    \begin{itemize}[<+->]
        \item Returning a string will simply show the string for you
        \item You can return a dictionary and it will dump JSON for you
        \item Returning file objects (or anything with \texttt{.read()}) will
            show the file contents for you
        \item You can even return an iterable or \texttt{yield} in your
            function and it will continually pull and show
    \end{itemize}
\end{frame}

\begin{frame}{Templates}
    \begin{onlyenv}<1>
    Bottle includes a templating engine called \texttt{SimpleTemplate}. You can
    use it like this:
    \end{onlyenv}

    {\small
    \inputminted{python3}{examples/template.py}
    }

    \begin{onlyenv}<2>
    \begin{block}{\texttt{views/hello\_template.tpl}}
        \small
        \inputminted{html}{examples/hello_template.tpl}
    \end{block}
    \end{onlyenv}
\end{frame}

\begin{frame}{Applications and Sub-applications}
    \small
    \inputminted{python3}{examples/apps.py}
\end{frame}

\begin{frame}{Deploying your application}
    Bottle has many ways to deploy, here's the two most common:

    \begin{enumerate}
        \item Using WSGI, simply create an \texttt{app.wsgi} file that imports
            your bottle app as \texttt{application}
        \item Using CGI (compatible with nearly any web server), just put
            \mintinline{python3}{run(server='cgi')} in your CGI script
    \end{enumerate}
\end{frame}

\begin{frame}[standout]
    \Huge
    Questions?
\end{frame}

\end{document}
